\chapter*{Convenciones}
\addcontentsline{toc}{chapter}{Convenciones}
\label{cap:notacion-pseudocodigo}

Este capítulo no se lee de corrido: es una referencia para volver
cuando aparezca una notación que no reconozcas. Si querés arrancar
con el contenido, saltá al capítulo~\ref{cap:threat-hunting}.

\section*{Notación matemática}

\begin{itemize}
    \item \(G = (\V, \E)\) — un grafo dirigido, donde \(\V\) es el
    conjunto de vértices y \(\E \subseteq \V \times \V\) el de aristas.
    Para grafos temporales, las aristas son ternas \((u, v, t)\) con
    \(t \in \T\); el \emph{multigrafo} permite que existan varias
    aristas entre el mismo par \((u, v)\) con distintos \(t\).
    \item \(|\V| = n\), \(|\E| = m\). Convención CLRS estándar.
    \item \(\Lc\) — alfabeto de etiquetas (en GraphHunter:
    \texttt{EntityType} para vértices, \texttt{RelationType} para
    aristas).
    \item \(\deg(v)\), \(\indeg(v)\), \(\outdeg(v)\) — grado, grado de
    entrada y de salida del vértice \(v\).
    \item \(N(v)\), \(N^{+}(v)\), \(N^{-}(v)\) — vecindad, vecindad
    saliente y vecindad entrante.
    \item \(\NLF(v, L)\) — \emph{Neighborhood Label Frequency}: número
    de vecinos salientes de \(v\) con etiqueta \(L\). Definida
    formalmente en el capítulo~\ref{cap:nlf-detalle}.
    \item \(\bits[v]\) — bitmap de \(|\Lc|\) bits asociado a \(v\),
    donde el bit \(L\) está prendido sii existe un camino válido de
    \(v\) a algún vértice de tipo \(L\). Definido en el
    capítulo~\ref{cap:khop-detalle}.
    \item \(\AGM(Q)\) y \(\rhostar(Q)\) — \emph{Atserias-Grohe-Marx
    bound} y \emph{fractional edge cover number} de la consulta \(Q\).
    Definidos y derivados en el capítulo~\ref{cap:wcoj}.
\end{itemize}

\section*{Notación del pseudocódigo}

Los algoritmos siguen convención CLRS, ligeramente adaptada:

\begin{itemize}
    \item Asignación: \(x \leftarrow y\). Igualdad: \(x = y\).
    \item Acceso a campos: \(v.\mathrm{campo}\) (notación punto, igual
    que un \emph{struct} en C/Rust).
    \item Bloques \texttt{for/while/if/else} delimitados por
    indentación más \texttt{end} explícito (igual que algorithm2e).
    \item \texttt{\textbf{return}} sale del procedimiento.
    \item \texttt{\textbf{emit}} agrega un resultado a la salida del
    enumerador (sin retornar).
    \item Conjuntos: \(\{x : P(x)\}\) (comprensión), \(A \cup B\),
    \(A \cap B\), \(A \setminus B\), \(\emptyset\).
    \item Operaciones \emph{bitwise}: \(a \mid b\) (\emph{or}),
    \(a \mathrel{|}\!= b\) (\emph{or-assign}), \(a \,\&\, b\)
    (\emph{and}). Cuando aparezcan, los explicamos en el momento.
\end{itemize}

\section*{Cajas y entornos visuales}

\begin{intuicion}
Los párrafos importantes en lenguaje plano arrancan con una caja
\emph{Intuición} verde clara. Su propósito es que entiendas el
\emph{qué} y el \emph{por qué} antes de leer el \emph{cómo} formal.
Si la prueba o el pseudocódigo se vuelven difíciles, releé la
intuición ---suele alcanzar para reconstruirlos.
\end{intuicion}

\begin{ejemplo}
Las cajas \emph{Ejemplo} (borde verde claro, fondo blanco)
contienen un ejemplo trabajado paso a paso. Suelen incluir un
diagrama TikZ y la traza de un algoritmo sobre una entrada chica.
\end{ejemplo}

\begin{enelcodigo}
Los párrafos \textbf{En el código} apuntan del concepto al fichero
real del repo donde está implementado. Por ejemplo:
\filepath{core/graph-engine/src/lftj.rs}.
\end{enelcodigo}

\begin{nota}
Las notas (sin numeración) son aclaraciones puntuales: una
diferencia con un libro de referencia, una pista para el lector
avanzado, una elección de diseño que vale la pena justificar.
\end{nota}

\begin{advertencia}
Las advertencias señalan errores frecuentes. Léelas con cuidado: por
ejemplo, ``no confundir \emph{path} con \emph{walk}'' es una trampa
clásica del threat hunting (¿se permite repetir vértices?).
\end{advertencia}

\section*{Teoremas, lemas, definiciones}

\begin{definition}[Plantilla]
Cada definición se enuncia con cuantificadores explícitos y la
referencia al objeto matemático preciso. La numeración es por
capítulo: \(\mathrm{Definici\acute{o}n}\,X.Y\) significa la \(Y\)-ésima
definición del capítulo \(X\).
\end{definition}

\begin{theorem}[Plantilla]
Cada teorema enuncia formalmente una propiedad y la prueba sigue
inmediatamente debajo. El símbolo \(\square\) cierra la prueba.
\end{theorem}

\begin{proof}
Esta es la prueba. \(\square\)
\end{proof}

Los \emph{Lemas} son resultados auxiliares; los \emph{Corolarios}
consecuencias inmediatas; las \emph{Proposiciones}, enunciados de
medio peso (más fuertes que un lema, menos que un teorema).

\section*{Ejercicios}

\begin{ejercicios}
Cada capítulo termina con 3--6 ejercicios. Tres tipos:

\begin{enumerate}
    \item \textbf{Warm-up}: confirma que entendiste las definiciones.
    Suelen ser de 1--3 minutos.
    \item \textbf{Derivaciones}: pedís \emph{trabajar} un ejemplo
    nuevo o probar una variante chica del teorema. 10--30 minutos.
    \item \textbf{Extensiones}: la solución requiere ir más allá del
    capítulo. Pueden ser problemas abiertos del motor; los abordamos
    en el capítulo siguiente o quedan como \emph{food for thought}.
\end{enumerate}

Los ejercicios marcados con \starred{} tienen solución en el
Apéndice~\ref{ap:soluciones}. Los demás quedan abiertos por diseño.
\end{ejercicios}

\section*{Referencias cruzadas}

\begin{itemize}
    \item Los \texttt{\textbackslash ref}s a capítulos son del estilo
    ``capítulo~\ref{cap:threat-hunting}''.
    \item Los \texttt{\textbackslash ref}s a apéndices son del estilo
    ``Apéndice~\ref{ap:notacion}''.
    \item Las figuras se referencian con ``figura~X.Y'' y los
    algoritmos con ``algoritmo~X.Y''.
    \item Las citas a literatura externa son nombre+año entre
    paréntesis: ``(Veldhuizen, 2014)''. Las referencias completas
    están en la bibliografía (Apéndice~\ref{ap:bibliografia}).
\end{itemize}

\cleardoublepage
